\documentclass[a4paper,10pt]{article} 
\usepackage{anysize}
\marginsize{2cm}{2cm}{1cm}{1cm}
%\textwidth 6.0in \textheight = 664pt
\usepackage{xltxtra}
\usepackage{graphicx}
\usepackage{subfig}
\usepackage{color}
\usepackage{xgreek}
\usepackage{framed}
\usepackage{float}
\usepackage{tkz-graph}
\setmainfont[Mapping=TeX-text]{CMU Serif}

%\setmainfont{Arial}
%\setsansfont{FreeSans}
%\setmonofont{FreeMono}
\begin{document}
\def\thesubsection {\alph{subsection}}
%\renewcommand{\labelenumi}{\roman{enumi})}
\renewcommand{\labelenumi}{\alph{enumi}}
\renewcommand{\labelenumii}{(\arabic{enumii})}


\begin{titlepage}
\begin{center}
\begin{figure}[t] 
     \includegraphics[scale=0.7]{title/ntua_logo}
\end{figure}
\begin{LARGE}\textbf{ΕΘΝΙΚΟ ΜΕΤΣΟΒΙΟ ΠΟΛΥΤΕΧΝΕΙΟ\\}\end{LARGE}
\vspace{5cm}
\begin{Large}
ΣΧΟΛΗ ΗΜ\&ΜΥ\\
Αλγόριθμοι και Πολυπλοκότητα\\
1\textsuperscript{η} Σειρά Γραπτών Ασκήσεων\\
Ακ. έτος 2010-2011\\
\end{Large}
\vspace{10cm}
\begin{tabular}{l r}
\Large{Λύρας Γρηγόρης}&
\large{Α.Μ.: 03109687}\\
\end{tabular}\\
\vspace{1cm}

\vfill
\large\today\\
\end{center}
\end{titlepage}


%{{{ Προβολή ταινιών
\section{Προβολή ταινιών} \setcounter{section}{1}

Κρατάμε για κάθε ταινία $M_i$ έναν ακέραιο $Delta_i$ τον οποίο αρχικοποιούμε
στο 0 καθώς και μια σημαία που δηλώνει αν αυτή η ταινία έχει επιλεγεί από
κάποιο συνδρομητή. Καθώς περνάμε τις προτιμήσεις του χρήστη $a$ σημειώνουμε
για την ταινία της επιλογής του τη σημαία της σε "ΕΠΙΛΕΓΜΕΝΗ" και αυξάνουμε
τον μετρητή της αν αυτή είναι σημειωμένη για προβολή Σαββάτου από τον χρήστη
ενώ μειώνουμε αν αυτή έχει επιλεγεί για Κυριακή. Αφού ολοκληρώσουμε για όλους
τους χρήστες. διαβάζουμε τον συντελεστή κάθε των ταινίας και αν αυτός είναι
θετικός την τοποθετούμε στο σύνολο του Σαββάτου αν είναι αρνητικός στο σύνολο
της Κυριακής. Αν είναι 0 και έχει σημαία "ΕΠΙΛΕΓΜΕΝΗ" δεν έχει σημασία πού θα
επιλέξουμε να την βάλουμε.  Αν η σημαία δεν είναι "ΕΠΙΛΕΓΜΕΝΗ" τότε αγνοούμε
αυτή την ταινία.
%}}}

%{{{ Μέτρηση συντομότερων μονοπατιών
\section{Μέτρηση Συντομότερων Μονοπατιών}
Ξεκινάμε από το άκρο s του γράφου. Θέτουμε τον κόμβο "υπό εξερεύνηση" και
ακολουθούμε όσες ακμές εξέρχονται από αυτόν. Θέτουμε τον κόμβο ως
"εξερευνημένο" και προχωρούμε στους νέους κόμβους. Αν με αυτή τη διαδικασία
φτάσουμε σε κόμβο "εξερευνημένο" είτε "υπό εξερεύνηση" τότε τον αγνούμε. Μόλις
φτάσουμε στο άκρο t ολοκληρώνουμε το τρέχον επίπεδι και σταματάμε τη
διαδικασία. Ο αριθμός των συντομότερων μονοπατιών είναι ο αριθμός των t που
εμφανίζεται στην τελευταία λίστα.

Η χρονική πολυπλοκότητα του αλγορίθμου είναι $\Theta(n)$
%}}}

%{{{ Ελάχιστο Συνδετικό Δέντρο Υπό Περιορισμούς
\section{Ελάχιστο Συνδετικό Δέντρο Υπό Περιορισμούς}
\subsection{}
\begin{figure}[ht]
	\centering
	\subfloat[Initial]{
		\begin{tikzpicture}

			  \Vertex[x=0,y=0,L=1]{1}
			  \Vertex[x=3,y=0,L=2]{2}
			  \Vertex[x=0,y=-3,L=3]{3}
			  \Vertex[x=3,y=-3,L=4]{4}

			  \Edge[label=2](1)(2)%
			  \Edge[label=8](1)(3)
			  \Edge[label=4](2)(3)
			  \Edge[label=3](3)(4)
	  \end{tikzpicture}
  }
  \hspace{1cm}
  \subfloat[MST]{
		\begin{tikzpicture}
			  \Vertex[x=0,y=0,L=1]{1}
			  \Vertex[x=3,y=0,L=2]{2}
			  \Vertex[x=0,y=-3,L=3]{3}
			  \Vertex[x=3,y=-3,L=4]{4}

			  \Edge[label=2](1)(2)%
			  \Edge[label=4](2)(3)
			  \Edge[label=3](3)(4)
	  \end{tikzpicture}
  }
  \hspace{1cm}
  \subfloat[L=\{2,4\}]{
		\begin{tikzpicture}
			  \Vertex[x=0,y=0,L=1]{1}
			  \Vertex[x=3,y=0,L=2]{2}
			  \Vertex[x=0,y=-3,L=3]{3}
			  \Vertex[x=3,y=-3,L=4]{4}

			  \Edge[label=2](1)(2)%
			  \Edge[label=8](1)(3)
			  \Edge[label=3](3)(4)
	  \end{tikzpicture}
  }
  \end{figure}

\subsection{}
	Αφαιρούμε τις κορυφές του L και βρίσκουμε το ελάχιστο συνδετικό δέντρο χωρίς
		αυτές. Τέλος προσθέτουμε τις κορυφές του L ως φύλλα στο δέντρο που
		φτιάξαμε χρησιμοποιώντας τις πιο "φτηνές" ακμές τους.

%}}}

%{{{ Μοναδικότητα Ελάχιστου Συνδετικού Δέντρου
\section{Μοναδικότητα Ελάχιστου Συνδετικού Δέντρου}
\begin{enumerate}
	\item
	\item
	\item
	\item
\end{enumerate}

%}}}

%{{{ Υπολογισμός ελάχιστου συνδετικού δέντρου με διαγραφή ακμών
\section{Υπολογισμός Ελάχιστου Συνδετικού Δέντρου με Διαγραφή Ακμών}
\begin{enumerate}
	\item
	\item
	\item
\end{enumerate}

%}}}
\end{document}
